\SourceLeaf{218}%
gleichen, indem wir das Grössenverhältniss ganz bei Seite setzen; als ähnlichliegend werden also zwei Darstellungen betrachtet, wenn von zwei aus Einem Punkte ausgehenden Linearelementen dem in der einen Darstellung rechts liegenden auch in der andern das rechts liegende entspricht: im entgegengesetzten Falle werden sie verkehrtliegende heissen. Bei der Ebne, von Nro. 2---7 wird immer die Seite, wo die positiven Werthe der dritten Coordinate liegen, als die obere betrachtet; bei der ersten und letzten Fläche hingegen ist die Unterscheidung der obern und untern Seite bloss von dem positiven oder negativen Werthe von \(\psi\) und \(\Psi\) abhängig, wie schon oben festgesetzt ist.

Hier ist nun zuvörderst klar, dass für jede Stelle der ersten Fläche, wo man bei ungeändertem \(x\) und \(y\) durch ein positives Increment von \(z\) auf deren obere Seite kommt, die Darstellung in 2 mit der in 1 ähnlichliegend sein wird; dies wird also offenbar überall zutreffen, wo \(h\) positiv ist; und das Gegentheil wird bei einem negativen \(h\) eintreten, wo die Darstellungen verkehrt liegend sein werden.

Auf dieselbe Weise werden die Darstellungen in 7 und 8 ähnlich liegend oder verkehrt liegend sein, jenachdem \(H\) positiv oder negativ ist.

Um die Darstellungen in 2 und 3 unter sich zu vergleichen, sei in der erstern \(\mathrm{d}s\) die Länge einer unendlich kleinen Linie von dem Punkte, dessen Coordinaten \(x,y\), zu einem andern, dessen Coordinaten \(x+\mathrm{d}x,y+\mathrm{d}y\) sind, und \(l\) dessen Neigung gegen die Abscissenlinie wachsend in dem Sinn, in welchem man von der Axe der \(x\) zu der Axe der \(y\) übergeht, also
\[
\mathrm{d}x=\mathrm{d}s\mathord{.}\cos l,\quad \mathrm{d}y=\mathrm{d}s\mathord{.}\sin l
\]
\noindent In der Darstellung 3 sei \(\mathrm{d}\sigma\) die Grösse der Linie, welche der \(\mathrm{d}s\) entspricht, und ihre Neigung zur Abscissenlinie, wie vorhin verstanden, \(\lambda\), so dass
\[
\mathrm{d}t=\mathrm{d}\sigma\mathord{.}\cos\lambda,\quad \mathrm{d}u=\mathrm{d}\sigma\mathord{.}\sin\lambda
\]
\noindent Man hat also, in den Bezeichnungen des 4. Artikels
\[
\begin{aligned}
\mathrm{d}s\mathord{.}\cos l&=\mathrm{d}\sigma(a\cos\lambda+a'\sin\lambda)\\
\mathrm{d}s\mathord{.}\sin l&=\mathrm{d}\sigma(b\cos\lambda+b'\sin\lambda)
\end{aligned}
\]
\noindent folglich
\[
\mathrm{tang}\,l=\frac{b\cos\lambda+b'\sin\lambda}{a\cos\lambda+a'\sin\lambda}
\]
